% preamble.tex - shared house preamble for /explain-projects documents.
% Every package here is VERIFIED present in this machine's TeX Live install.
% Do not add packages without running: kpsewhich <pkg>.sty
% Notably ABSENT (do not \usepackage them, the build will fail):
%   algorithm2e, algorithm, algorithmic, algpseudocode  (texlive-science is not installed)
%   -> use the \begin{pseudocode} tcolorbox environment defined below instead.
%   lmodern, charter, mathptmx, times, mathpazo, courier, newtx*, tgtermes
%   -> their .sty files may resolve but the FONT FILES are absent, so the build
%      dies with "I can't find file `bchr8t'" or similar. Verified 2026-07-16 by
%      test-compiling each one. Only Computer Modern (the default) and helvet
%      actually render on this box. Do not "upgrade" the font stack without
%      test-compiling first: kpsewhich lying about a font is the exact trap here.
%
% To get a nicer font stack + algorithm2e, install:
%   sudo apt install texlive-fonts-recommended texlive-science
% Until then, this preamble deliberately uses the default font.
%
% Usage from a document:
%   \documentclass[11pt,a4paper]{article}
%   % preamble.tex - shared house preamble for /explain-projects documents.
% Every package here is VERIFIED present in this machine's TeX Live install.
% Do not add packages without running: kpsewhich <pkg>.sty
% Notably ABSENT (do not \usepackage them, the build will fail):
%   algorithm2e, algorithm, algorithmic, algpseudocode  (texlive-science is not installed)
%   -> use the \begin{pseudocode} tcolorbox environment defined below instead.
%   lmodern, charter, mathptmx, times, mathpazo, courier, newtx*, tgtermes
%   -> their .sty files may resolve but the FONT FILES are absent, so the build
%      dies with "I can't find file `bchr8t'" or similar. Verified 2026-07-16 by
%      test-compiling each one. Only Computer Modern (the default) and helvet
%      actually render on this box. Do not "upgrade" the font stack without
%      test-compiling first: kpsewhich lying about a font is the exact trap here.
%
% To get a nicer font stack + algorithm2e, install:
%   sudo apt install texlive-fonts-recommended texlive-science
% Until then, this preamble deliberately uses the default font.
%
% Usage from a document:
%   \documentclass[11pt,a4paper]{article}
%   % preamble.tex - shared house preamble for /explain-projects documents.
% Every package here is VERIFIED present in this machine's TeX Live install.
% Do not add packages without running: kpsewhich <pkg>.sty
% Notably ABSENT (do not \usepackage them, the build will fail):
%   algorithm2e, algorithm, algorithmic, algpseudocode  (texlive-science is not installed)
%   -> use the \begin{pseudocode} tcolorbox environment defined below instead.
%   lmodern, charter, mathptmx, times, mathpazo, courier, newtx*, tgtermes
%   -> their .sty files may resolve but the FONT FILES are absent, so the build
%      dies with "I can't find file `bchr8t'" or similar. Verified 2026-07-16 by
%      test-compiling each one. Only Computer Modern (the default) and helvet
%      actually render on this box. Do not "upgrade" the font stack without
%      test-compiling first: kpsewhich lying about a font is the exact trap here.
%
% To get a nicer font stack + algorithm2e, install:
%   sudo apt install texlive-fonts-recommended texlive-science
% Until then, this preamble deliberately uses the default font.
%
% Usage from a document:
%   \documentclass[11pt,a4paper]{article}
%   % preamble.tex - shared house preamble for /explain-projects documents.
% Every package here is VERIFIED present in this machine's TeX Live install.
% Do not add packages without running: kpsewhich <pkg>.sty
% Notably ABSENT (do not \usepackage them, the build will fail):
%   algorithm2e, algorithm, algorithmic, algpseudocode  (texlive-science is not installed)
%   -> use the \begin{pseudocode} tcolorbox environment defined below instead.
%   lmodern, charter, mathptmx, times, mathpazo, courier, newtx*, tgtermes
%   -> their .sty files may resolve but the FONT FILES are absent, so the build
%      dies with "I can't find file `bchr8t'" or similar. Verified 2026-07-16 by
%      test-compiling each one. Only Computer Modern (the default) and helvet
%      actually render on this box. Do not "upgrade" the font stack without
%      test-compiling first: kpsewhich lying about a font is the exact trap here.
%
% To get a nicer font stack + algorithm2e, install:
%   sudo apt install texlive-fonts-recommended texlive-science
% Until then, this preamble deliberately uses the default font.
%
% Usage from a document:
%   \documentclass[11pt,a4paper]{article}
%   \input{preamble}          % this file is copied next to the .tex

\usepackage[T1]{fontenc}
\usepackage[utf8]{inputenc}
% Font: Computer Modern (default). See the note above before changing this.
\usepackage[a4paper,margin=2.4cm,headheight=14pt]{geometry}
% microtype WITHOUT font expansion. This is load-bearing, do not "restore" it:
% cm-super is not installed, so T1 Computer Modern resolves to BITMAP fonts, and
% pdfTeX's auto-expansion needs scalable ones. With expansion on, any document
% long enough to reach a second page dies with
%   "pdfTeX error (font expansion): auto expansion is only possible with scalable fonts"
% Protrusion (the bigger visual win) still works and stays on.
% `sudo apt install cm-super` would lift this; until then, leave expansion off.
\usepackage[protrusion=true,expansion=false]{microtype}
\usepackage{graphicx}
\usepackage{xcolor}
\usepackage{booktabs}
\usepackage{enumitem}
\usepackage{fancyhdr}
\usepackage{titlesec}
\usepackage{tcolorbox}
\usepackage{listings}
\usepackage{float}
\usepackage{caption}
\usepackage{tikz}
\usepackage{forest}
\usepackage{pgfplots}
\pgfplotsset{compat=1.18}
\usetikzlibrary{arrows.meta,positioning,shapes.geometric,fit,backgrounds,calc,decorations.pathreplacing}
\tcbuselibrary{skins,breakable}
\usepackage[hidelinks,bookmarks=true,bookmarksnumbered=true]{hyperref}

% ===== palette ===========================================================
\definecolor{inkblue}{HTML}{1F3A5F}
\definecolor{accent}{HTML}{2E7DAF}
\definecolor{softgrey}{HTML}{F4F6F8}
\definecolor{linegrey}{HTML}{C8D0D8}
\definecolor{codebg}{HTML}{F7F7F4}
\definecolor{warnred}{HTML}{A33B3B}
\definecolor{okgreen}{HTML}{2E7D53}

% ===== headings ==========================================================
\titleformat{\section}{\Large\bfseries\color{inkblue}}{\thesection}{0.7em}{}
\titleformat{\subsection}{\large\bfseries\color{inkblue}}{\thesubsection}{0.6em}{}
\titleformat{\subsubsection}{\normalsize\bfseries\color{accent}}{\thesubsubsection}{0.5em}{}

\pagestyle{fancy}
\fancyhf{}
\fancyhead[L]{\small\color{inkblue}\nouppercase{\leftmark}}
\fancyhead[R]{\small\color{inkblue}\thepage}
\renewcommand{\headrulewidth}{0.4pt}
\renewcommand{\headrule}{\hbox to\headwidth{\color{linegrey}\leaders\hrule height \headrulewidth\hfill}}

% ===== code ==============================================================
\lstdefinestyle{house}{
  backgroundcolor=\color{codebg},
  basicstyle=\ttfamily\footnotesize,
  keywordstyle=\color{inkblue}\bfseries,
  commentstyle=\color{gray}\itshape,
  stringstyle=\color{okgreen},
  numbers=left, numberstyle=\tiny\color{gray}, numbersep=8pt,
  frame=single, rulecolor=\color{linegrey},
  breaklines=true, breakatwhitespace=false,
  showstringspaces=false, tabsize=2,
  captionpos=b, xleftmargin=14pt, framexleftmargin=14pt,
  columns=fullflexible, keepspaces=true,
  upquote=true  % render ' and " as straight quotes, not curly
}
\lstset{style=house}

% ===== callout boxes =====================================================
% \begin{keyidea}{Title} ... \end{keyidea}   - a concept a beginner must grasp
\newtcolorbox{keyidea}[1]{
  enhanced, breakable, colback=softgrey, colframe=accent,
  boxrule=0.4pt, arc=2pt, left=8pt, right=8pt, top=6pt, bottom=6pt,
  title={\bfseries #1}, coltitle=white, colbacktitle=accent, fonttitle=\small
}

% \begin{jargon}{Term} definition \end{jargon}  - define a term at first use
\newtcolorbox{jargon}[1]{
  enhanced, breakable, colback=white, colframe=linegrey,
  boxrule=0.4pt, arc=2pt, left=8pt, right=8pt, top=6pt, bottom=6pt,
  title={\bfseries Jargon: #1}, coltitle=inkblue, colbacktitle=softgrey, fonttitle=\small
}

% \begin{honest}{Title} ... \end{honest}  - brutally honest finding (rule 18)
\newtcolorbox{honest}[1]{
  enhanced, breakable, colback=white, colframe=warnred,
  boxrule=0.6pt, arc=2pt, left=8pt, right=8pt, top=6pt, bottom=6pt,
  title={\bfseries #1}, coltitle=white, colbacktitle=warnred, fonttitle=\small
}

% \begin{provenance}{path/to/file} ... \end{provenance}
% Where a file came from: hand-written, generated by a tool, scaffolded by a
% framework, copied or adapted from an upstream project, ported from an older
% module. State the EVIDENCE (a header comment, a git log entry, an upstream
% URL, a generator command) and label the confidence. If no evidence exists,
% say so in the box. Never guess a lineage into existence.
\newtcolorbox{provenance}[1]{
  enhanced, breakable, colback=white, colframe=okgreen,
  boxrule=0.4pt, arc=2pt, left=8pt, right=8pt, top=6pt, bottom=6pt,
  title={\bfseries Where this came from: #1}, coltitle=white,
  colbacktitle=okgreen, fonttitle=\small
}

% \begin{pseudocode}{Caption} ... \end{pseudocode}
% Stands in for algorithm2e, which is NOT installed. Use \Step / \Indent inside.
\newtcolorbox{pseudocode}[1]{
  enhanced, breakable, colback=codebg, colframe=inkblue,
  boxrule=0.4pt, arc=2pt, left=10pt, right=8pt, top=6pt, bottom=6pt,
  fontupper=\ttfamily\small,
  title={\bfseries\rmfamily #1}, coltitle=white, colbacktitle=inkblue, fonttitle=\small
}
\newcommand{\Step}[1]{\par\noindent #1}
\newcommand{\Indent}[1]{\par\noindent\hspace*{2em}#1}

% ===== tikz house styles =================================================
\tikzset{
  box/.style   = {rectangle, rounded corners=2pt, draw=inkblue, fill=softgrey,
                  text=inkblue, align=center, inner sep=6pt, minimum height=9mm,
                  font=\small},
  extbox/.style= {rectangle, rounded corners=2pt, draw=linegrey, fill=white,
                  text=black!70, align=center, inner sep=6pt, minimum height=9mm,
                  font=\small, dashed},
  store/.style = {cylinder, shape border rotate=90, aspect=0.25, draw=accent,
                  fill=white, text=inkblue, align=center, inner sep=5pt,
                  font=\small},
  dec/.style   = {diamond, aspect=1.6, draw=accent, fill=white, text=inkblue,
                  align=center, inner sep=2pt, font=\scriptsize},
  flow/.style  = {-{Stealth[length=2.4mm]}, draw=inkblue, thick},
  dflow/.style = {-{Stealth[length=2.4mm]}, draw=accent, thick, dashed},
  % align=center is load-bearing: without it, a \\ inside an edge label throws
  % the misleading "perhaps a missing \item" error.
  lbl/.style   = {font=\scriptsize\itshape, text=black!65, fill=white,
                  inner sep=2pt, align=center}
}

% forest default for directory trees
\forestset{
  dirtree/.style={
    for tree={
      font=\ttfamily\small, grow'=0, child anchor=west, parent anchor=south,
      anchor=west, calign=first, inner sep=2pt,
      edge path={\noexpand\path [draw, \forestoption{edge}]
        (!u.south west) +(9pt,0) |- (.child anchor) \forestoption{edge label};},
      before typesetting nodes={if n=1{insert before={[,phantom]}}{}},
      fit=band, before computing xy={l=15pt}
    }
  }
}

% ===== misc ==============================================================
\captionsetup{font=small, labelfont={bf,color=inkblue}}
\setlist[itemize]{leftmargin=1.2em, itemsep=2pt, topsep=3pt}
\setlist[enumerate]{leftmargin=1.4em, itemsep=2pt, topsep=3pt}
\newcommand{\file}[1]{\texttt{\small #1}}
\newcommand{\code}[1]{\texttt{\small #1}}
\setcounter{tocdepth}{2}
\setcounter{secnumdepth}{3}
          % this file is copied next to the .tex

\usepackage[T1]{fontenc}
\usepackage[utf8]{inputenc}
% Font: Computer Modern (default). See the note above before changing this.
\usepackage[a4paper,margin=2.4cm,headheight=14pt]{geometry}
% microtype WITHOUT font expansion. This is load-bearing, do not "restore" it:
% cm-super is not installed, so T1 Computer Modern resolves to BITMAP fonts, and
% pdfTeX's auto-expansion needs scalable ones. With expansion on, any document
% long enough to reach a second page dies with
%   "pdfTeX error (font expansion): auto expansion is only possible with scalable fonts"
% Protrusion (the bigger visual win) still works and stays on.
% `sudo apt install cm-super` would lift this; until then, leave expansion off.
\usepackage[protrusion=true,expansion=false]{microtype}
\usepackage{graphicx}
\usepackage{xcolor}
\usepackage{booktabs}
\usepackage{enumitem}
\usepackage{fancyhdr}
\usepackage{titlesec}
\usepackage{tcolorbox}
\usepackage{listings}
\usepackage{float}
\usepackage{caption}
\usepackage{tikz}
\usepackage{forest}
\usepackage{pgfplots}
\pgfplotsset{compat=1.18}
\usetikzlibrary{arrows.meta,positioning,shapes.geometric,fit,backgrounds,calc,decorations.pathreplacing}
\tcbuselibrary{skins,breakable}
\usepackage[hidelinks,bookmarks=true,bookmarksnumbered=true]{hyperref}

% ===== palette ===========================================================
\definecolor{inkblue}{HTML}{1F3A5F}
\definecolor{accent}{HTML}{2E7DAF}
\definecolor{softgrey}{HTML}{F4F6F8}
\definecolor{linegrey}{HTML}{C8D0D8}
\definecolor{codebg}{HTML}{F7F7F4}
\definecolor{warnred}{HTML}{A33B3B}
\definecolor{okgreen}{HTML}{2E7D53}

% ===== headings ==========================================================
\titleformat{\section}{\Large\bfseries\color{inkblue}}{\thesection}{0.7em}{}
\titleformat{\subsection}{\large\bfseries\color{inkblue}}{\thesubsection}{0.6em}{}
\titleformat{\subsubsection}{\normalsize\bfseries\color{accent}}{\thesubsubsection}{0.5em}{}

\pagestyle{fancy}
\fancyhf{}
\fancyhead[L]{\small\color{inkblue}\nouppercase{\leftmark}}
\fancyhead[R]{\small\color{inkblue}\thepage}
\renewcommand{\headrulewidth}{0.4pt}
\renewcommand{\headrule}{\hbox to\headwidth{\color{linegrey}\leaders\hrule height \headrulewidth\hfill}}

% ===== code ==============================================================
\lstdefinestyle{house}{
  backgroundcolor=\color{codebg},
  basicstyle=\ttfamily\footnotesize,
  keywordstyle=\color{inkblue}\bfseries,
  commentstyle=\color{gray}\itshape,
  stringstyle=\color{okgreen},
  numbers=left, numberstyle=\tiny\color{gray}, numbersep=8pt,
  frame=single, rulecolor=\color{linegrey},
  breaklines=true, breakatwhitespace=false,
  showstringspaces=false, tabsize=2,
  captionpos=b, xleftmargin=14pt, framexleftmargin=14pt,
  columns=fullflexible, keepspaces=true,
  upquote=true  % render ' and " as straight quotes, not curly
}
\lstset{style=house}

% ===== callout boxes =====================================================
% \begin{keyidea}{Title} ... \end{keyidea}   - a concept a beginner must grasp
\newtcolorbox{keyidea}[1]{
  enhanced, breakable, colback=softgrey, colframe=accent,
  boxrule=0.4pt, arc=2pt, left=8pt, right=8pt, top=6pt, bottom=6pt,
  title={\bfseries #1}, coltitle=white, colbacktitle=accent, fonttitle=\small
}

% \begin{jargon}{Term} definition \end{jargon}  - define a term at first use
\newtcolorbox{jargon}[1]{
  enhanced, breakable, colback=white, colframe=linegrey,
  boxrule=0.4pt, arc=2pt, left=8pt, right=8pt, top=6pt, bottom=6pt,
  title={\bfseries Jargon: #1}, coltitle=inkblue, colbacktitle=softgrey, fonttitle=\small
}

% \begin{honest}{Title} ... \end{honest}  - brutally honest finding (rule 18)
\newtcolorbox{honest}[1]{
  enhanced, breakable, colback=white, colframe=warnred,
  boxrule=0.6pt, arc=2pt, left=8pt, right=8pt, top=6pt, bottom=6pt,
  title={\bfseries #1}, coltitle=white, colbacktitle=warnred, fonttitle=\small
}

% \begin{provenance}{path/to/file} ... \end{provenance}
% Where a file came from: hand-written, generated by a tool, scaffolded by a
% framework, copied or adapted from an upstream project, ported from an older
% module. State the EVIDENCE (a header comment, a git log entry, an upstream
% URL, a generator command) and label the confidence. If no evidence exists,
% say so in the box. Never guess a lineage into existence.
\newtcolorbox{provenance}[1]{
  enhanced, breakable, colback=white, colframe=okgreen,
  boxrule=0.4pt, arc=2pt, left=8pt, right=8pt, top=6pt, bottom=6pt,
  title={\bfseries Where this came from: #1}, coltitle=white,
  colbacktitle=okgreen, fonttitle=\small
}

% \begin{pseudocode}{Caption} ... \end{pseudocode}
% Stands in for algorithm2e, which is NOT installed. Use \Step / \Indent inside.
\newtcolorbox{pseudocode}[1]{
  enhanced, breakable, colback=codebg, colframe=inkblue,
  boxrule=0.4pt, arc=2pt, left=10pt, right=8pt, top=6pt, bottom=6pt,
  fontupper=\ttfamily\small,
  title={\bfseries\rmfamily #1}, coltitle=white, colbacktitle=inkblue, fonttitle=\small
}
\newcommand{\Step}[1]{\par\noindent #1}
\newcommand{\Indent}[1]{\par\noindent\hspace*{2em}#1}

% ===== tikz house styles =================================================
\tikzset{
  box/.style   = {rectangle, rounded corners=2pt, draw=inkblue, fill=softgrey,
                  text=inkblue, align=center, inner sep=6pt, minimum height=9mm,
                  font=\small},
  extbox/.style= {rectangle, rounded corners=2pt, draw=linegrey, fill=white,
                  text=black!70, align=center, inner sep=6pt, minimum height=9mm,
                  font=\small, dashed},
  store/.style = {cylinder, shape border rotate=90, aspect=0.25, draw=accent,
                  fill=white, text=inkblue, align=center, inner sep=5pt,
                  font=\small},
  dec/.style   = {diamond, aspect=1.6, draw=accent, fill=white, text=inkblue,
                  align=center, inner sep=2pt, font=\scriptsize},
  flow/.style  = {-{Stealth[length=2.4mm]}, draw=inkblue, thick},
  dflow/.style = {-{Stealth[length=2.4mm]}, draw=accent, thick, dashed},
  % align=center is load-bearing: without it, a \\ inside an edge label throws
  % the misleading "perhaps a missing \item" error.
  lbl/.style   = {font=\scriptsize\itshape, text=black!65, fill=white,
                  inner sep=2pt, align=center}
}

% forest default for directory trees
\forestset{
  dirtree/.style={
    for tree={
      font=\ttfamily\small, grow'=0, child anchor=west, parent anchor=south,
      anchor=west, calign=first, inner sep=2pt,
      edge path={\noexpand\path [draw, \forestoption{edge}]
        (!u.south west) +(9pt,0) |- (.child anchor) \forestoption{edge label};},
      before typesetting nodes={if n=1{insert before={[,phantom]}}{}},
      fit=band, before computing xy={l=15pt}
    }
  }
}

% ===== misc ==============================================================
\captionsetup{font=small, labelfont={bf,color=inkblue}}
\setlist[itemize]{leftmargin=1.2em, itemsep=2pt, topsep=3pt}
\setlist[enumerate]{leftmargin=1.4em, itemsep=2pt, topsep=3pt}
\newcommand{\file}[1]{\texttt{\small #1}}
\newcommand{\code}[1]{\texttt{\small #1}}
\setcounter{tocdepth}{2}
\setcounter{secnumdepth}{3}
          % this file is copied next to the .tex

\usepackage[T1]{fontenc}
\usepackage[utf8]{inputenc}
% Font: Computer Modern (default). See the note above before changing this.
\usepackage[a4paper,margin=2.4cm,headheight=14pt]{geometry}
% microtype WITHOUT font expansion. This is load-bearing, do not "restore" it:
% cm-super is not installed, so T1 Computer Modern resolves to BITMAP fonts, and
% pdfTeX's auto-expansion needs scalable ones. With expansion on, any document
% long enough to reach a second page dies with
%   "pdfTeX error (font expansion): auto expansion is only possible with scalable fonts"
% Protrusion (the bigger visual win) still works and stays on.
% `sudo apt install cm-super` would lift this; until then, leave expansion off.
\usepackage[protrusion=true,expansion=false]{microtype}
\usepackage{graphicx}
\usepackage{xcolor}
\usepackage{booktabs}
\usepackage{enumitem}
\usepackage{fancyhdr}
\usepackage{titlesec}
\usepackage{tcolorbox}
\usepackage{listings}
\usepackage{float}
\usepackage{caption}
\usepackage{tikz}
\usepackage{forest}
\usepackage{pgfplots}
\pgfplotsset{compat=1.18}
\usetikzlibrary{arrows.meta,positioning,shapes.geometric,fit,backgrounds,calc,decorations.pathreplacing}
\tcbuselibrary{skins,breakable}
\usepackage[hidelinks,bookmarks=true,bookmarksnumbered=true]{hyperref}

% ===== palette ===========================================================
\definecolor{inkblue}{HTML}{1F3A5F}
\definecolor{accent}{HTML}{2E7DAF}
\definecolor{softgrey}{HTML}{F4F6F8}
\definecolor{linegrey}{HTML}{C8D0D8}
\definecolor{codebg}{HTML}{F7F7F4}
\definecolor{warnred}{HTML}{A33B3B}
\definecolor{okgreen}{HTML}{2E7D53}

% ===== headings ==========================================================
\titleformat{\section}{\Large\bfseries\color{inkblue}}{\thesection}{0.7em}{}
\titleformat{\subsection}{\large\bfseries\color{inkblue}}{\thesubsection}{0.6em}{}
\titleformat{\subsubsection}{\normalsize\bfseries\color{accent}}{\thesubsubsection}{0.5em}{}

\pagestyle{fancy}
\fancyhf{}
\fancyhead[L]{\small\color{inkblue}\nouppercase{\leftmark}}
\fancyhead[R]{\small\color{inkblue}\thepage}
\renewcommand{\headrulewidth}{0.4pt}
\renewcommand{\headrule}{\hbox to\headwidth{\color{linegrey}\leaders\hrule height \headrulewidth\hfill}}

% ===== code ==============================================================
\lstdefinestyle{house}{
  backgroundcolor=\color{codebg},
  basicstyle=\ttfamily\footnotesize,
  keywordstyle=\color{inkblue}\bfseries,
  commentstyle=\color{gray}\itshape,
  stringstyle=\color{okgreen},
  numbers=left, numberstyle=\tiny\color{gray}, numbersep=8pt,
  frame=single, rulecolor=\color{linegrey},
  breaklines=true, breakatwhitespace=false,
  showstringspaces=false, tabsize=2,
  captionpos=b, xleftmargin=14pt, framexleftmargin=14pt,
  columns=fullflexible, keepspaces=true,
  upquote=true  % render ' and " as straight quotes, not curly
}
\lstset{style=house}

% ===== callout boxes =====================================================
% \begin{keyidea}{Title} ... \end{keyidea}   - a concept a beginner must grasp
\newtcolorbox{keyidea}[1]{
  enhanced, breakable, colback=softgrey, colframe=accent,
  boxrule=0.4pt, arc=2pt, left=8pt, right=8pt, top=6pt, bottom=6pt,
  title={\bfseries #1}, coltitle=white, colbacktitle=accent, fonttitle=\small
}

% \begin{jargon}{Term} definition \end{jargon}  - define a term at first use
\newtcolorbox{jargon}[1]{
  enhanced, breakable, colback=white, colframe=linegrey,
  boxrule=0.4pt, arc=2pt, left=8pt, right=8pt, top=6pt, bottom=6pt,
  title={\bfseries Jargon: #1}, coltitle=inkblue, colbacktitle=softgrey, fonttitle=\small
}

% \begin{honest}{Title} ... \end{honest}  - brutally honest finding (rule 18)
\newtcolorbox{honest}[1]{
  enhanced, breakable, colback=white, colframe=warnred,
  boxrule=0.6pt, arc=2pt, left=8pt, right=8pt, top=6pt, bottom=6pt,
  title={\bfseries #1}, coltitle=white, colbacktitle=warnred, fonttitle=\small
}

% \begin{provenance}{path/to/file} ... \end{provenance}
% Where a file came from: hand-written, generated by a tool, scaffolded by a
% framework, copied or adapted from an upstream project, ported from an older
% module. State the EVIDENCE (a header comment, a git log entry, an upstream
% URL, a generator command) and label the confidence. If no evidence exists,
% say so in the box. Never guess a lineage into existence.
\newtcolorbox{provenance}[1]{
  enhanced, breakable, colback=white, colframe=okgreen,
  boxrule=0.4pt, arc=2pt, left=8pt, right=8pt, top=6pt, bottom=6pt,
  title={\bfseries Where this came from: #1}, coltitle=white,
  colbacktitle=okgreen, fonttitle=\small
}

% \begin{pseudocode}{Caption} ... \end{pseudocode}
% Stands in for algorithm2e, which is NOT installed. Use \Step / \Indent inside.
\newtcolorbox{pseudocode}[1]{
  enhanced, breakable, colback=codebg, colframe=inkblue,
  boxrule=0.4pt, arc=2pt, left=10pt, right=8pt, top=6pt, bottom=6pt,
  fontupper=\ttfamily\small,
  title={\bfseries\rmfamily #1}, coltitle=white, colbacktitle=inkblue, fonttitle=\small
}
\newcommand{\Step}[1]{\par\noindent #1}
\newcommand{\Indent}[1]{\par\noindent\hspace*{2em}#1}

% ===== tikz house styles =================================================
\tikzset{
  box/.style   = {rectangle, rounded corners=2pt, draw=inkblue, fill=softgrey,
                  text=inkblue, align=center, inner sep=6pt, minimum height=9mm,
                  font=\small},
  extbox/.style= {rectangle, rounded corners=2pt, draw=linegrey, fill=white,
                  text=black!70, align=center, inner sep=6pt, minimum height=9mm,
                  font=\small, dashed},
  store/.style = {cylinder, shape border rotate=90, aspect=0.25, draw=accent,
                  fill=white, text=inkblue, align=center, inner sep=5pt,
                  font=\small},
  dec/.style   = {diamond, aspect=1.6, draw=accent, fill=white, text=inkblue,
                  align=center, inner sep=2pt, font=\scriptsize},
  flow/.style  = {-{Stealth[length=2.4mm]}, draw=inkblue, thick},
  dflow/.style = {-{Stealth[length=2.4mm]}, draw=accent, thick, dashed},
  % align=center is load-bearing: without it, a \\ inside an edge label throws
  % the misleading "perhaps a missing \item" error.
  lbl/.style   = {font=\scriptsize\itshape, text=black!65, fill=white,
                  inner sep=2pt, align=center}
}

% forest default for directory trees
\forestset{
  dirtree/.style={
    for tree={
      font=\ttfamily\small, grow'=0, child anchor=west, parent anchor=south,
      anchor=west, calign=first, inner sep=2pt,
      edge path={\noexpand\path [draw, \forestoption{edge}]
        (!u.south west) +(9pt,0) |- (.child anchor) \forestoption{edge label};},
      before typesetting nodes={if n=1{insert before={[,phantom]}}{}},
      fit=band, before computing xy={l=15pt}
    }
  }
}

% ===== misc ==============================================================
\captionsetup{font=small, labelfont={bf,color=inkblue}}
\setlist[itemize]{leftmargin=1.2em, itemsep=2pt, topsep=3pt}
\setlist[enumerate]{leftmargin=1.4em, itemsep=2pt, topsep=3pt}
\newcommand{\file}[1]{\texttt{\small #1}}
\newcommand{\code}[1]{\texttt{\small #1}}
\setcounter{tocdepth}{2}
\setcounter{secnumdepth}{3}
          % this file is copied next to the .tex

\usepackage[T1]{fontenc}
\usepackage[utf8]{inputenc}
% Font: Computer Modern (default). See the note above before changing this.
\usepackage[a4paper,margin=2.4cm,headheight=14pt]{geometry}
% microtype WITHOUT font expansion. This is load-bearing, do not "restore" it:
% cm-super is not installed, so T1 Computer Modern resolves to BITMAP fonts, and
% pdfTeX's auto-expansion needs scalable ones. With expansion on, any document
% long enough to reach a second page dies with
%   "pdfTeX error (font expansion): auto expansion is only possible with scalable fonts"
% Protrusion (the bigger visual win) still works and stays on.
% `sudo apt install cm-super` would lift this; until then, leave expansion off.
\usepackage[protrusion=true,expansion=false]{microtype}
\usepackage{graphicx}
\usepackage{xcolor}
\usepackage{booktabs}
\usepackage{enumitem}
\usepackage{fancyhdr}
\usepackage{titlesec}
\usepackage{tcolorbox}
\usepackage{listings}
\usepackage{float}
\usepackage{caption}
\usepackage{tikz}
\usepackage{forest}
\usepackage{pgfplots}
\pgfplotsset{compat=1.18}
\usetikzlibrary{arrows.meta,positioning,shapes.geometric,fit,backgrounds,calc,decorations.pathreplacing}
\tcbuselibrary{skins,breakable}
\usepackage[hidelinks,bookmarks=true,bookmarksnumbered=true]{hyperref}

% ===== palette ===========================================================
\definecolor{inkblue}{HTML}{1F3A5F}
\definecolor{accent}{HTML}{2E7DAF}
\definecolor{softgrey}{HTML}{F4F6F8}
\definecolor{linegrey}{HTML}{C8D0D8}
\definecolor{codebg}{HTML}{F7F7F4}
\definecolor{warnred}{HTML}{A33B3B}
\definecolor{okgreen}{HTML}{2E7D53}

% ===== headings ==========================================================
\titleformat{\section}{\Large\bfseries\color{inkblue}}{\thesection}{0.7em}{}
\titleformat{\subsection}{\large\bfseries\color{inkblue}}{\thesubsection}{0.6em}{}
\titleformat{\subsubsection}{\normalsize\bfseries\color{accent}}{\thesubsubsection}{0.5em}{}

\pagestyle{fancy}
\fancyhf{}
\fancyhead[L]{\small\color{inkblue}\nouppercase{\leftmark}}
\fancyhead[R]{\small\color{inkblue}\thepage}
\renewcommand{\headrulewidth}{0.4pt}
\renewcommand{\headrule}{\hbox to\headwidth{\color{linegrey}\leaders\hrule height \headrulewidth\hfill}}

% ===== code ==============================================================
\lstdefinestyle{house}{
  backgroundcolor=\color{codebg},
  basicstyle=\ttfamily\footnotesize,
  keywordstyle=\color{inkblue}\bfseries,
  commentstyle=\color{gray}\itshape,
  stringstyle=\color{okgreen},
  numbers=left, numberstyle=\tiny\color{gray}, numbersep=8pt,
  frame=single, rulecolor=\color{linegrey},
  breaklines=true, breakatwhitespace=false,
  showstringspaces=false, tabsize=2,
  captionpos=b, xleftmargin=14pt, framexleftmargin=14pt,
  columns=fullflexible, keepspaces=true,
  upquote=true  % render ' and " as straight quotes, not curly
}
\lstset{style=house}

% ===== callout boxes =====================================================
% \begin{keyidea}{Title} ... \end{keyidea}   - a concept a beginner must grasp
\newtcolorbox{keyidea}[1]{
  enhanced, breakable, colback=softgrey, colframe=accent,
  boxrule=0.4pt, arc=2pt, left=8pt, right=8pt, top=6pt, bottom=6pt,
  title={\bfseries #1}, coltitle=white, colbacktitle=accent, fonttitle=\small
}

% \begin{jargon}{Term} definition \end{jargon}  - define a term at first use
\newtcolorbox{jargon}[1]{
  enhanced, breakable, colback=white, colframe=linegrey,
  boxrule=0.4pt, arc=2pt, left=8pt, right=8pt, top=6pt, bottom=6pt,
  title={\bfseries Jargon: #1}, coltitle=inkblue, colbacktitle=softgrey, fonttitle=\small
}

% \begin{honest}{Title} ... \end{honest}  - brutally honest finding (rule 18)
\newtcolorbox{honest}[1]{
  enhanced, breakable, colback=white, colframe=warnred,
  boxrule=0.6pt, arc=2pt, left=8pt, right=8pt, top=6pt, bottom=6pt,
  title={\bfseries #1}, coltitle=white, colbacktitle=warnred, fonttitle=\small
}

% \begin{provenance}{path/to/file} ... \end{provenance}
% Where a file came from: hand-written, generated by a tool, scaffolded by a
% framework, copied or adapted from an upstream project, ported from an older
% module. State the EVIDENCE (a header comment, a git log entry, an upstream
% URL, a generator command) and label the confidence. If no evidence exists,
% say so in the box. Never guess a lineage into existence.
\newtcolorbox{provenance}[1]{
  enhanced, breakable, colback=white, colframe=okgreen,
  boxrule=0.4pt, arc=2pt, left=8pt, right=8pt, top=6pt, bottom=6pt,
  title={\bfseries Where this came from: #1}, coltitle=white,
  colbacktitle=okgreen, fonttitle=\small
}

% \begin{pseudocode}{Caption} ... \end{pseudocode}
% Stands in for algorithm2e, which is NOT installed. Use \Step / \Indent inside.
\newtcolorbox{pseudocode}[1]{
  enhanced, breakable, colback=codebg, colframe=inkblue,
  boxrule=0.4pt, arc=2pt, left=10pt, right=8pt, top=6pt, bottom=6pt,
  fontupper=\ttfamily\small,
  title={\bfseries\rmfamily #1}, coltitle=white, colbacktitle=inkblue, fonttitle=\small
}
\newcommand{\Step}[1]{\par\noindent #1}
\newcommand{\Indent}[1]{\par\noindent\hspace*{2em}#1}

% ===== tikz house styles =================================================
\tikzset{
  box/.style   = {rectangle, rounded corners=2pt, draw=inkblue, fill=softgrey,
                  text=inkblue, align=center, inner sep=6pt, minimum height=9mm,
                  font=\small},
  extbox/.style= {rectangle, rounded corners=2pt, draw=linegrey, fill=white,
                  text=black!70, align=center, inner sep=6pt, minimum height=9mm,
                  font=\small, dashed},
  store/.style = {cylinder, shape border rotate=90, aspect=0.25, draw=accent,
                  fill=white, text=inkblue, align=center, inner sep=5pt,
                  font=\small},
  dec/.style   = {diamond, aspect=1.6, draw=accent, fill=white, text=inkblue,
                  align=center, inner sep=2pt, font=\scriptsize},
  flow/.style  = {-{Stealth[length=2.4mm]}, draw=inkblue, thick},
  dflow/.style = {-{Stealth[length=2.4mm]}, draw=accent, thick, dashed},
  % align=center is load-bearing: without it, a \\ inside an edge label throws
  % the misleading "perhaps a missing \item" error.
  lbl/.style   = {font=\scriptsize\itshape, text=black!65, fill=white,
                  inner sep=2pt, align=center}
}

% forest default for directory trees
\forestset{
  dirtree/.style={
    for tree={
      font=\ttfamily\small, grow'=0, child anchor=west, parent anchor=south,
      anchor=west, calign=first, inner sep=2pt,
      edge path={\noexpand\path [draw, \forestoption{edge}]
        (!u.south west) +(9pt,0) |- (.child anchor) \forestoption{edge label};},
      before typesetting nodes={if n=1{insert before={[,phantom]}}{}},
      fit=band, before computing xy={l=15pt}
    }
  }
}

% ===== misc ==============================================================
\captionsetup{font=small, labelfont={bf,color=inkblue}}
\setlist[itemize]{leftmargin=1.2em, itemsep=2pt, topsep=3pt}
\setlist[enumerate]{leftmargin=1.4em, itemsep=2pt, topsep=3pt}
\newcommand{\file}[1]{\texttt{\small #1}}
\newcommand{\code}[1]{\texttt{\small #1}}
\setcounter{tocdepth}{2}
\setcounter{secnumdepth}{3}
